\documentclass[11pt]{article}
\usepackage[margin=1in]{geometry}
\usepackage{amsmath,amssymb,mathtools,bm}
\usepackage{booktabs}
\usepackage{hyperref}

\title{First-Step-Contraction MPC With a Bounded Frozen-\texttt{dhat} Steady-State Target}
\author{}
\date{\today}

\begin{document}
\maketitle

\section{Purpose}

This report documents the bounded frozen-\texttt{dhat} target method that is now available beside the existing refined Step A selector in the first-step-contraction MPC workflow.

The intent is narrow:
\begin{itemize}
\item keep the maintained upstream offset-free MPC and first-step-contraction logic unchanged,
\item replace only the steady-state target package used as the Lyapunov center,
\item preserve the old refined-selector path in the codebase and in the original notebook.
\end{itemize}

The new notebook uses the same upstream controller, the same replacement logic, and the same debug-export path. Only the way the steady-state package $(x_s,u_s,d_s,y_s,r_s)$ is computed changes.

\section{Scaled Deviation Coordinates}

All quantities below are in the same scaled deviation coordinates already used by the identified linear model, the observer, and the MPC controller. Let
\begin{align}
u_k^{\mathrm{dev}} &= \mathrm{scale}(u_k^{\mathrm{phys}}) - u_{ss}^{\mathrm{scaled}}, \\
y_k^{\mathrm{dev}} &= \mathrm{scale}(y_k^{\mathrm{phys}}) - y_{ss}^{\mathrm{scaled}}.
\end{align}
The steady-state target calculation is performed entirely in these deviation coordinates. Conversion back to physical units is done only for plotting or reporting.

\section{Augmented Offset-Free Model}

The first-step-contraction notebook uses the same Rawlings-style output-disturbance augmentation as the maintained offset-free MPC workflow:
\begin{align}
z_k &=
\begin{bmatrix}
x_k \\
d_k
\end{bmatrix}, \\
z_{k+1} &=
\underbrace{\begin{bmatrix}
A & B_d \\
0 & I
\end{bmatrix}}_{A_{\mathrm{aug}}}
z_k
+
\underbrace{\begin{bmatrix}
B \\
0
\end{bmatrix}}_{B_{\mathrm{aug}}}
u_k, \\
y_k &=
\underbrace{\begin{bmatrix}
C & C_d
\end{bmatrix}}_{C_{\mathrm{aug}}}
z_k.
\end{align}

In the notebook configuration used for the new method:
\begin{align}
C_d = I,
\end{align}
and the augmented state is ordered as
\begin{align}
z_k = \begin{bmatrix} x_k \\ d_k \end{bmatrix},
\end{align}
with the last $n_y$ augmented states representing output disturbances.

\section{Observer Split and Frozen Disturbance Assumption}

At each control step the observer provides the augmented estimate
\begin{align}
\hat z_k =
\begin{bmatrix}
\hat x_k \\
\hat d_k
\end{bmatrix}.
\end{align}

The bounded frozen-\texttt{dhat} target method uses the simplest steady-state disturbance policy:
\begin{align}
d_s(k) = \hat d_k.
\end{align}
So the disturbance target is \emph{not} optimized. It is frozen at the current observer estimate.

This differs from the refined Step A selector, where the target package is produced by a separate convex selector problem with its own smoothing and anchoring penalties. Here there is no secondary selector cost shaping the steady package. The disturbance estimate itself defines the output offset.

\section{Frozen-\texttt{dhat} Steady-State Equations}

Because the target computation is based on the unaugmented plant model with output disturbance entering through the measurement equation, the steady-state conditions are
\begin{align}
x_s &= A x_s + B u_s, \\
y_s &= C x_s + d_s,
\end{align}
with
\begin{align}
d_s = \hat d_k.
\end{align}

If the requested setpoint at step $k$ is $y_{sp,k}$, then the exact offset-free target must satisfy
\begin{align}
y_s = y_{sp,k}.
\end{align}
Substituting $d_s = \hat d_k$ gives
\begin{align}
C x_s = y_{sp,k} - \hat d_k.
\end{align}

Define the right-hand side
\begin{align}
r_k := y_{sp,k} - \hat d_k.
\end{align}
Then the frozen-\texttt{dhat} steady-state equations become
\begin{align}
(I-A)x_s - B u_s &= 0, \\
C x_s &= r_k.
\end{align}

\section{Stacked Exact Solve}

Introduce the stacked unknown
\begin{align}
w_s :=
\begin{bmatrix}
x_s \\
u_s
\end{bmatrix}.
\end{align}
Then the steady-state equations can be written as
\begin{align}
M w_s = b_k,
\end{align}
with
\begin{align}
M &=
\begin{bmatrix}
I-A & -B \\
C   & 0
\end{bmatrix}, \\
b_k &=
\begin{bmatrix}
0 \\
r_k
\end{bmatrix}.
\end{align}

The exact solution is obtained from this linear system whenever the system is square, full-rank, and sufficiently well-conditioned. In code this is the \texttt{stacked\_exact} branch.

\section{Reduced Exact Solve}

If $(I-A)$ is safely invertible, then
\begin{align}
x_s = (I-A)^{-1} B u_s.
\end{align}
Substituting this into the output equation yields the reduced map
\begin{align}
G := C(I-A)^{-1}B,
\end{align}
and the reduced steady-state equation
\begin{align}
G u_s = r_k.
\end{align}

If $G$ is square, full-rank, and well-conditioned, the exact target is
\begin{align}
u_s = G^{-1} r_k, \qquad
x_s = (I-A)^{-1} B u_s.
\end{align}

The code uses the following exact-solve priority:
\begin{enumerate}
\item \texttt{stacked\_exact} if the stacked system is square, full-rank, and well-conditioned,
\item \texttt{reduced\_exact} if the reduced map is available and well-conditioned,
\item least-squares fallbacks only if the exact modes are not numerically trustworthy.
\end{enumerate}

\section{Why a Box-Constrained Fallback Is Needed}

The exact frozen-\texttt{dhat} target may violate the input box
\begin{align}
u_{\min} \le u_s \le u_{\max}.
\end{align}
When this happens, the exact target is not an admissible steady-state center for a constrained controller. It is still diagnostically useful, but it is not a proper feasible target.

That is why the second stage solves a simple bounded optimization problem.

\section{Bounded Least-Squares Fallback}

The simplest bounded version of the target problem is
\begin{align}
\min_{x_s,u_s} \quad
\left\|
\begin{bmatrix}
I-A & -B \\
C   & 0
\end{bmatrix}
\begin{bmatrix}
x_s \\
u_s
\end{bmatrix}
-
\begin{bmatrix}
0 \\
r_k
\end{bmatrix}
\right\|_2^2
\qquad
\text{subject to }
u_{\min} \le u_s \le u_{\max}.
\end{align}

If the reduced map is available, the same problem can be written more compactly as
\begin{align}
\min_{u_s} \quad \|G u_s - r_k\|_2^2
\qquad
\text{subject to }
u_{\min} \le u_s \le u_{\max},
\end{align}
and then
\begin{align}
x_s = (I-A)^{-1}B u_s.
\end{align}

This is exactly the ``simplest optimization version'' described in the design discussion.

\section{Projection Interpretation}

Assume the reduced map $G$ is invertible and define the unbounded exact target
\begin{align}
u_s^{\mathrm{exact}} = G^{-1} r_k.
\end{align}
Then the bounded objective becomes
\begin{align}
\|G u - r_k\|_2^2
=
\|G(u-u_s^{\mathrm{exact}})\|_2^2
=
(u-u_s^{\mathrm{exact}})^\top G^\top G (u-u_s^{\mathrm{exact}}).
\end{align}

So the bounded fallback is the projection of the exact target onto the admissible input box under the metric
\begin{align}
H = G^\top G.
\end{align}
This explains an important practical point:
\begin{itemize}
\item if the exact target lies inside the box, the bounded and exact solutions coincide,
\item if the exact target lies outside the box, the bounded solution often lands on a box face or corner.
\end{itemize}

\section{KKT Conditions and Active Bounds}

For the reduced bounded problem define
\begin{align}
J(u) = \frac{1}{2}\|G u - r_k\|_2^2.
\end{align}
Its gradient is
\begin{align}
\nabla J(u) = G^\top(G u - r_k).
\end{align}

The Karush--Kuhn--Tucker conditions at an optimum $u^\star$ are
\begin{align}
u_{\min} \le u^\star \le u_{\max}, \\
G^\top(G u^\star - r_k) + \lambda^+ - \lambda^- = 0, \\
\lambda^+ \ge 0, \qquad \lambda^- \ge 0, \\
\lambda_i^+(u_i^\star - u_{\max,i}) = 0, \\
\lambda_i^-(u_{\min,i} - u_i^\star) = 0.
\end{align}

These conditions explain why a bound-active steady-state target is expected in disturbance cases: if the exact target is outside the box, the optimal constrained solution is often characterized by nonzero multipliers on active bounds.

\section{Returned Steady-State Package}

The new target helper returns a package compatible with the first-step-contraction workflow:
\begin{align}
x_s, \qquad
u_s, \qquad
d_s = \hat d_k, \qquad
y_s = Cx_s + d_s, \qquad
r_s = y_s.
\end{align}

It also returns
\begin{align}
x_s^{\mathrm{aug}} =
\begin{bmatrix}
x_s \\
d_s
\end{bmatrix},
\end{align}
plus diagnostics such as
\begin{align}
\|r_s-y_{sp,k}\|_\infty, \quad
\|(I-A)x_s-Bu_s\|_\infty, \quad
\text{bound violation}, \quad
\text{exact-within-bounds flag}, \quad
\text{bounded residual norm}.
\end{align}

The key semantic choices are:
\begin{itemize}
\item \texttt{d\_s\_frozen = True},
\item \texttt{d\_s\_optimized = False},
\item \texttt{d\_s\_minus\_dhat\_inf = 0}.
\end{itemize}

\section{How the New Target Enters the First-Step Controller}

The first-step-contraction controller still follows the maintained architecture:
\begin{enumerate}
\item compute a target package,
\item build the effective target using current target or last valid target,
\item solve the ordinary upstream offset-free MPC candidate,
\item check first-step Lyapunov contraction relative to the effective target,
\item if needed, solve a constrained replacement MPC problem with one additional first-step Lyapunov inequality.
\end{enumerate}

Only Step 1 changes.

\section{Upstream Offset-Free MPC Candidate}

With the augmented model
\begin{align}
z_{j+1|k} = A_{\mathrm{aug}} z_{j|k} + B_{\mathrm{aug}} u_{c(j)|k},
\qquad
c(j)=\min(j,N_C-1),
\end{align}
the predicted output is
\begin{align}
y_{j|k} = C_{\mathrm{aug}} z_{j|k}.
\end{align}

The upstream controller still solves the baseline offset-free MPC objective
\begin{align}
J_{\mathrm{MPC}}(U)
=
\sum_{j=1}^{N_P}
\|y_{j|k} - y_{\mathrm{track},k}\|_{Q_y}^2
+
\sum_{j=0}^{N_C-1}
\|\Delta u_{j|k}\|_{R_{\Delta u}}^2.
\end{align}

In the new notebook,
\begin{align}
y_{\mathrm{track},k} = y_{sp,k},
\end{align}
because both
\texttt{mpc\_target\_policy} and
\texttt{tracking\_target\_policy}
remain \texttt{raw\_setpoint}.

So the new steady-state method affects the Lyapunov center, not the nominal tracking objective of the upstream MPC.

\section{Lyapunov Function}

Let $P_x \succ 0$ be the physical-state Lyapunov matrix used by the maintained first-step-contraction code. Define the physical-state error relative to the effective target:
\begin{align}
e_{x,k} = \hat x_k - x_s^{\mathrm{eff}}.
\end{align}

The Lyapunov function is
\begin{align}
V_k = e_{x,k}^\top P_x e_{x,k}.
\end{align}
The contraction bound is
\begin{align}
V_{\mathrm{bound},k} = \rho_{\mathrm{lyap}} V_k + \varepsilon_{\mathrm{lyap}},
\qquad
0 < \rho_{\mathrm{lyap}} < 1.
\end{align}

\section{First-Step Contraction Inequality}

For a candidate first move $u_{0|k}$, the predicted first physical-state step is
\begin{align}
\hat x_{k+1|k}^{(1)} = A \hat x_k + B u_{0|k} + B_d \hat d_k,
\end{align}
using the augmented model that the controller itself predicts with.

The corresponding predicted Lyapunov value is
\begin{align}
V_{k+1|k}^{(1)}
=
(\hat x_{k+1|k}^{(1)} - x_s^{\mathrm{eff}})^\top
P_x
(\hat x_{k+1|k}^{(1)} - x_s^{\mathrm{eff}}).
\end{align}

The maintained first-step condition is
\begin{align}
V_{k+1|k}^{(1)} \le \rho_{\mathrm{lyap}} V_k + \varepsilon_{\mathrm{lyap}}.
\end{align}

If the upstream candidate already satisfies this inequality, it is accepted unchanged. If not, the controller triggers the replacement solve.

\section{Replacement Solve}

The replacement controller keeps the same MPC objective and the same move/input constraints as the upstream problem, but adds the single extra hard first-step constraint
\begin{align}
V_{k+1|k}^{(1)} \le \rho_{\mathrm{lyap}} V_k + \varepsilon_{\mathrm{lyap}}.
\end{align}

So the modified optimization is still
\begin{align}
\min_U \quad J_{\mathrm{MPC}}(U)
\end{align}
subject to
\begin{align}
& z_{j+1|k} = A_{\mathrm{aug}} z_{j|k} + B_{\mathrm{aug}} u_{c(j)|k}, \\
& u_{\min} \le u_{j|k} \le u_{\max}, \\
& V_{k+1|k}^{(1)} \le \rho_{\mathrm{lyap}} V_k + \varepsilon_{\mathrm{lyap}}.
\end{align}

There is still
\begin{itemize}
\item no QCQP projection stage,
\item no post-hoc safety filter,
\item no change to the upstream stage cost,
\item no change to the candidate reuse and backup logic.
\end{itemize}

\section{Nominal Versus Disturbance Interpretation}

The nominal case and the disturbance case behave differently even if the target solver is implemented correctly.

\subsection{Nominal Case}

If
\begin{align}
\hat d_k \approx 0,
\end{align}
then
\begin{align}
r_k = y_{sp,k} - \hat d_k \approx y_{sp,k}.
\end{align}
If the loop settles, both the plant and the frozen-\texttt{dhat} target tend toward the same nominal equilibrium. In that case it is common to see
\begin{align}
u_{\mathrm{applied},k} \to u_s(k).
\end{align}

\subsection{Disturbance Case}

In the disturbance case, the target is driven by
\begin{align}
r_k = y_{sp,k} - \hat d_k.
\end{align}
Even when $r_k$ becomes small, the reduced map may amplify that small residual:
\begin{align}
u_s = G^{-1} r_k.
\end{align}
If $G^{-1}$ has large entries, then small changes in $r_k$ can still produce visible motion in $u_s$.

Also, the controller does \emph{not} track $u_s$ directly. The upstream MPC still tracks the raw output setpoint, while the bounded frozen-\texttt{dhat} target is used only as the Lyapunov center. Therefore there is no reason to expect
\begin{align}
u_{\mathrm{applied},k} = u_s(k)
\end{align}
in the disturbance case, even near apparent steady output behavior.

\section{Algorithmic Summary}

At each control step $k$ the new target-generation branch executes:
\begin{enumerate}
\item Split $\hat z_k$ into $\hat x_k$ and $\hat d_k$.
\item Form $r_k = y_{sp,k} - \hat d_k$.
\item Solve the exact frozen-\texttt{dhat} steady-state equations using the unaugmented $A,B,C$.
\item Check whether the exact input target lies inside the input box.
\item If yes, return the exact target package.
\item If no, solve the bounded least-squares fallback and return the bounded target package.
\item Hand that package to the existing effective-target, upstream MPC, and first-step-contraction logic.
\end{enumerate}

\section{Code-Level Mapping}

The implementation is split as follows:
\begin{itemize}
\item \texttt{Lyapunov/frozen\_dhat\_target.py}: computes the bounded frozen-\texttt{dhat} target package.
\item \texttt{Simulation/run\_mpc\_first\_step\_contraction.py}: adds a new \texttt{target\_generation\_mode} branch.
\item \texttt{LyapunovFirstStepContractionMPC\_BoundedFrozenDhat.ipynb}: uses the new branch.
\item \texttt{LyapunovFirstStepContractionMPC.ipynb}: remains on the refined-selector path.
\end{itemize}

\section{Conclusion}

The bounded frozen-\texttt{dhat} method is intentionally simpler than the refined Step A selector. It produces the Lyapunov center by solving the frozen-disturbance steady-state equations directly, then projecting onto the admissible input box when the exact target is infeasible. This gives a target that is easier to interpret mathematically:
\begin{itemize}
\item exact equilibrium if the box permits it,
\item closest feasible equilibrium surrogate in the least-squares sense otherwise.
\end{itemize}

The rest of the first-step-contraction MPC architecture remains unchanged.

\end{document}
